\paragraph{逐维度可利用性分析。}
下表将攻击方式与对应检查明确配对。“有边界”表示可利用但受约束；分类适用于当前测试域。任务包和研究数据保留具体示例与缓解措施。
\begin{center}\small\begin{tabular}{p{.20\linewidth}p{.17\linewidth}p{.52\linewidth}}\toprule
维度 & 分类 & 攻击方式；对应检查 \\ \midrule
程序与素材身份 & 难以利用 & 非法程序、替换素材；受限解释与身份检查。 \\
物理与接触 & 有边界 & 利用接触容差；实际网格与未知状态拒绝通过。 \\
操作目标 & 有边界 & 无碰撞但未完成操作；显式目标验收。 \\
软约束与诊断 & 有边界 & 虚报不可满足；联合见证与声明核对。 \\
度量重建 & 有边界 & 轮廓相似但中心错误；联合度量验收。 \\
轨迹与事件 & 有边界 & 只满足端点；连续检查与实际释放。 \\
编辑与保留 & 有边界 & 不修改以降低代价；目标与保留验收。 \\
检查器一致性 & 有边界 & 依赖错误通过；几何与诊断分别审计。 \\
视觉质量 & 软性判断 & 遮挡或外观；辅助角色与覆盖率。 \\
目标访问 & 有边界 & 读取隐藏文件；受限公开接口。 \\ \bottomrule
\end{tabular}\end{center}
